\section{Generalized Wasserstein Distances}
\label{sec-extensions}
\label{sec-generalized-wasserstein-distances}
\subsection{Unbalanced OT}
\label{sec-unbalanced}
\paragraph{Relaxed formulation.}
For nonnegative measures $(\al,\be) \in \Mm_+(\Xx) \times \Mm_+(\Yy)$, a generic relaxed formulation is

\eql{\label{eq-unbalanced-primal}
\UW_c(\al,\be)= \uinf{\pi \in \Mm_+(\Xx \times \Yy)} \int_{\Xx \times \Yy} c(x,y) \d\pi(x,y)
+ \Dd_{\psi_1}(\pi_1|\al) + \Dd_{\psi_2}(\pi_2|\be),
}
where $\psi_1,\psi_2$ are convex entropy functions. Exact conservation $(\pi_1,\pi_2)=(\al,\be)$ is replaced by a cost for changing the marginals.

\(\UW_{c,\tau}(\alpha,\beta) = \inf_{\pi\geq0} \int c\d\pi + \tau\Dd_{\bar\psi_1}(\pi_1|\alpha) + \tau\Dd_{\bar\psi_2}(\pi_2|\beta).\)
These figures should be read as pictures of a single relaxed plan $\pi$: the same nonnegative measure determines both the transported coupling and its two relaxed marginals. The small-$\tau$ result below formalizes the opposite regime, where transport becomes negligible compared with local mass variation.

\begin{prop}[Small-transport-scale limit for marginal penalties]\label{prop-unbalanced-small-scale-limit}
Assume that $\alpha,\beta$ are finite measures on a compact metric space $\Xx$, that $c$ is continuous, $c\geq0$, and $c(x,y)=0$ if and only if $x=y$. Assume also that the marginal divergences are nonnegative, weak-* lower semicontinuous, and have weak-* compact sublevel sets on $\Mm_+(\Xx)$. Then
\[
\lim_{\tau\downarrow0}\frac1\tau\UW_{c,\tau}(\alpha,\beta)
=
\inf_{\rho\in\Mm_+(\Xx)}
\Dd_{\bar\psi_1}(\rho|\alpha)
+
\Dd_{\bar\psi_2}(\rho|\beta).
\]
The right-hand side is the infimal gluing divergence obtained by matching the two measures through a common zero-transport marginal $\rho$. In the dominated case, if $\alpha=a\lambda$, $\beta=b\lambda$, and the minimizing common marginal is absolutely continuous, $\rho=r\lambda$, this divergence decouples pointwise as
\[
\int \mathfrak m_{\bar\psi_1,\bar\psi_2}(a(x),b(x))\d\lambda(x),
\qquad
\mathfrak m_{\bar\psi_1,\bar\psi_2}(a,b)
\eqdef
\inf_{r\geq0} a\,\bar\psi_1(r/a)+b\,\bar\psi_2(r/b),
\]
with the usual recession conventions when $a=0$ or $b=0$. For superlinear entropies, and in particular for KL, finite energy forces this dominated form. Thus, when $\bar\psi_1=\bar\psi_2$ is the KL entropy,
\(\inf_{\rho\in\Mm_+(\Xx)} \KL(\rho|\alpha)+\KL(\rho|\beta) = \int (\sqrt{a}-\sqrt{b})^2\d\lambda.\)
Thus the KL marginal relaxation contains the squared Hellinger distance as its local mass-variation limit.
\end{prop}
\begin{proof}
For the upper bound, restrict to diagonal plans $\pi=(\Id,\Id)_\sharp\rho$, whose transport cost is zero and whose two marginals are both $\rho$.

For the lower bound, let $\tau_n\downarrow0$ and let $\pi_n$ be almost minimizing plans with bounded scaled values $\tau_n^{-1}\UW_{c,\tau_n}(\alpha,\beta)$. Since the divergences are nonnegative, $\int c\d\pi_n=O(\tau_n)$, hence $\int c\d\pi_n\to0$. The bounded scaled values also put the two marginals in compact divergence sublevel sets. Since a coupling has the same total mass as each marginal, the couplings are tight on $\Xx\times\Xx$. Up to subsequences, $\pi_n\rightharpoonup\pi_0$. Lower semicontinuity of the transport cost yields $\int c\d\pi_0=0$, so $\pi_0$ is concentrated on the diagonal. Its two marginals are therefore equal to a common measure $\rho$. Lower semicontinuity of the marginal divergences gives
\(\liminf_n\frac1{\tau_n}\UW_{c,\tau_n}(\alpha,\beta) \geq \Dd_{\bar\psi_1}(\rho|\alpha)+\Dd_{\bar\psi_2}(\rho|\beta),\)
and optimizing over $\rho$ gives the lower bound.

In the dominated case, writing $\rho=r\lambda$ gives \(\Dd_{\bar\psi_1}(\rho|\alpha)+\Dd_{\bar\psi_2}(\rho|\beta) = \int a\,\bar\psi_1(r/a)+b\,\bar\psi_2(r/b) \d\lambda,\)
so the minimization over $\rho$ decouples into the scalar envelope $\mathfrak m_{\bar\psi_1,\bar\psi_2}$. For KL, no singular part is admissible when $\alpha$ and $\beta$ are dominated by $\lambda$. The pointwise objective is $r\log(r/a)-r+a+r\log(r/b)-r+b$. Its optimality condition is $\log(r/a)+\log(r/b)=0$, hence $r=\sqrt{ab}$, and the minimum is $a+b-2\sqrt{ab}=(\sqrt a-\sqrt b)^2$.
\end{proof}
\begin{prop}[Dual of unbalanced optimal transport]\label{prop-dual-unbalanced-ot}
Under the usual Fenchel--Rockafellar qualification assumptions, one has equality between~\eqref{eq-unbalanced-primal} and
\(\UW_c(\al,\be) = \usup{f\oplus g\leq c} - \Dd_{\psi_1}^*(-f|\al) - \Dd_{\psi_2}^*(-g|\be).\)
\end{prop}
\begin{proof}
Use the variational formula~\eqref{eq-legendre} for the dual of a divergence and introduce the marginal variables through continuous potentials:
\(\inf_{\pi\geq0}\sup_{f,g} \int c\d\pi + \int -f\d\pi_1 + \int -g\d\pi_2 - \Dd_{\psi_1}^*(-f|\al) - \Dd_{\psi_2}^*(-g|\be).\)
Exchanging the infimum and the supremum gives \(\sup_{f,g} - \Dd_{\psi_1}^*(-f|\al) - \Dd_{\psi_2}^*(-g|\be) + \inf_{\pi\geq0}\int\big(c-(f\oplus g)\big)\d\pi.\)
The last infimum is $0$ when $f\oplus g\leq c$ and $-\infty$ otherwise, which gives the displayed dual.
\end{proof}
\paragraph{Reverse and homogeneous formulations.}
\begin{align*}
&\int c(x,y)\d\pi(x,y)
+\Dd_{\psi_1}(\pi_1|\al)+\Dd_{\psi_2}(\pi_2|\be) \\
&\qquad =
\int
\left(
c(x,y)
+
\psi_1\!\left(\frac{\d\pi_1}{\d\al}(x)\right)\frac{\d\al}{\d\pi_1}(x)
+
\psi_2\!\left(\frac{\d\pi_2}{\d\be}(y)\right)\frac{\d\be}{\d\pi_2}(y)
\right)
\d\pi(x,y).
\end{align*}
\eql{\label{eq-unbalanced-reverse-local-cost}
L_c(r,s) \eqdef c + r\psi_1(1/r)+s\psi_2(1/s),
}
with the usual recession convention at $r=0$ or $s=0$. If $\al=F\pi_1+\al^\perp$ and $\be=G\pi_2+\be^\perp$ are the Lebesgue decompositions of the reference marginals with respect to the transported marginals, then the reverse formulation reads

\(\UW_c(\al,\be) = \inf_{\pi\geq0} \int L_{c(x,y)}(F(x),G(y))\d\pi(x,y) + \psi_1(0)\al^\perp(\Xx) + \psi_2(0)\be^\perp(\Yy).\)
\eql{\label{eq-unbalanced-homogeneous-local-cost}
H_c(r,s) \eqdef \inf_{\theta>0}\theta L_c(r/\theta,s/\theta),
}
\eql{\label{eq-homogeneous}
\HW_c(\al,\be)
=
\inf_{\pi\geq0}
\int H_{c(x,y)}(F(x),G(y))\d\pi(x,y)
+
\psi_1(0)\al^\perp(\Xx)
+
\psi_2(0)\be^\perp(\Yy).
}
\eql{\label{eq-homogeneous-semicoupling}
\begin{aligned}
\HW_c(\al,\be)
=
\inf_{\substack{\lambda\in\Mm_+(\Xx\times\Yy)\\ u,v\geq0}}
\Big\{
&\int_{\Xx\times\Yy}
H_{c(x,y)}\big(u(x,y),v(x,y)\big)\d\lambda(x,y)
\\[-.15em]
&\text{subject to}\quad (\mathrm{p}_1)_\#(u\lambda)=\al,\qquad
(\mathrm{p}_2)_\#(v\lambda)=\be
\Big\}.
\end{aligned}
}
\begin{prop}[Homogenization does not change the unbalanced cost]\label{prop-homogeneous-unbalanced}
Assume that the entropy functions are proper, lower semicontinuous and convex, with the usual recession extensions at zero. Then
\(\HW_c(\al,\be)=\UW_c(\al,\be).\)
\end{prop}
\begin{proof}
Since $H_c(r,s)\leq L_c(r,s)$ by choosing the scale equal to one, every competitor in the reverse formulation gives a homogeneous competitor of no larger cost. Hence $\HW_c\leq\UW_c$.

For the converse, use the semi-coupling representation~\eqref{eq-homogeneous-semicoupling} and fix a feasible triple $(\lambda,u,v)$. For a positive measurable scale $\theta(x,y)$, set $\tilde\pi=\theta\lambda$. Disintegrate $\lambda$ with respect to its first spatial marginal. If $\bar u(x)$ and $\bar\theta(x)$ denote the corresponding conditional averages of $u$ and $\theta$, then $\alpha=\bar u\,(\mathrm p_1)_\sharp\lambda$ and $\tilde\pi_1=\bar\theta\,(\mathrm p_1)_\sharp\lambda$. Convexity of the perspective $(a,m)\mapsto a\psi_1(m/a)$ and conditional Jensen give
\(\Dd_{\psi_1}(\tilde\pi_1|\alpha) \leq \int u\,\psi_1(\theta/u)\d\lambda.\)
The same argument on the second marginal yields the analogous bound with $(v,\psi_2)$. Therefore
\[
\int c\d\tilde\pi
+\Dd_{\psi_1}(\tilde\pi_1|\alpha)
+\Dd_{\psi_2}(\tilde\pi_2|\beta)
\leq
\int\big[c\theta+u\psi_1(\theta/u)+v\psi_2(\theta/v)\big]\d\lambda.
\]
The pointwise infimum of the integrand over $\theta>0$ is precisely $H_c(u,v)$. A measurable $\eta$-minimizing choice of $\theta$ exists because this is a lower-semicontinuous normal integrand; the recession conventions cover $u=0$ or $v=0$. Letting $\eta\downarrow0$ and then minimizing over $(\lambda,u,v)$ gives $\UW_c\leq\HW_c$.
\end{proof}
\paragraph{Conic lifting.}
Assume now that $\Xx=\Yy$ and $\psi_1=\psi_2=\psi$. The last formulation lifts the problem to the cone space $\mathfrak C[\Xx]\eqdef(\Xx\times\RR_+)/\sim$, where all points $(x,0)$ are identified at the apex. For an exponent $p\geq1$, define

\(\De((x,r),(y,s))\eqdef H_{c(x,y)}(r^p,s^p)^{1/p}.\)
\begin{itemize}
\item $\Dd_\psi=\KL$, $p=2$, and $c(x,y)=-\log\cos^2(d(x,y)\wedge\pi/2)$ give the Hellinger--Kantorovich or Wasserstein--Fisher--Rao cone metric
\(\De((x,r),(y,s))^2=r^2+s^2-2rs\cos(d(x,y)\wedge\pi/2).\)
\item $\Dd_\psi=\KL$, $p=2$, and $c(x,y)=d(x,y)^2$ give the Gaussian Hellinger formula
\(\De((x,r),(y,s))^2=r^2+s^2-2rs e^{-d(x,y)^2/2}.\)
It is a cone metric whenever the Gaussian kernel $k(x,y)=e^{-d(x,y)^2/2}$ is positive definite. This holds, for example, when $(\Xx,d)$ is a subset of a Hilbert space: if $k(x,y)=\langle\Phi(x),\Phi(y)\rangle$ with $\|\Phi(x)\|=1$, then the displayed quantity is $\|r\Phi(x)-s\Phi(y)\|^2$. For a general metric space, positive definiteness of this kernel is an additional hypothesis.
\item $\Dd_\psi=\TV$, $p=1$, and $c(x,y)=d(x,y)$ give the partial-transport cone cost
\(\De((x,r),(y,s))=r+s-(r\wedge s)(2-d(x,y))_+.\)
\end{itemize}
For a finite measure $\eta$ on the cone, define its weighted base projection $\mathsf P_p\eta\in\Mm_+(\Xx)$ by

\(\int_{\Xx}\varphi(x)\d(\mathsf P_p\eta)(x) = \int_{\mathfrak C[\Xx]}\varphi(x)r^p\d\eta(x,r)\)
for every bounded Borel $\varphi$. This is well defined at the apex because the weight $r^p$ vanishes there.

\[
\CW(\al,\be)
=
\inf_{\gamma\in\Mm_+(\mathfrak C[\Xx]^2)}
\left\{
\int \De((x,r),(y,s))^p\d\gamma
\;; \;
\mathsf P_p\gamma_1=\al,\quad
\mathsf P_p\gamma_2=\be
\right\}.
\]
\begin{thm}[Cone formulation of unbalanced OT]\label{thm-cone-unbalanced-ot}
Assume that $\Xx$ is a compact metric space, that $(r,s)\mapsto H_{c(x,y)}(r,s)$ is convex for every $(x,y)$, and that $(x,y,r,s)\mapsto H_{c(x,y)}(r,s)$ is lower semicontinuous. Then
\(\UW_c(\alpha,\beta)=\HW_c(\alpha,\beta)=\CW(\alpha,\beta).\)
If, in addition, $\De$ is a continuous distance on $\mathfrak C[\Xx]$, then $\CW^{1/p}$ is a distance between nonnegative finite measures.
\end{thm}
\begin{proof}
The equality $\UW=\HW$ is Proposition~\ref{prop-homogeneous-unbalanced}.

Let $(\lambda,u,v)$ be feasible in~\eqref{eq-homogeneous-semicoupling}. Define the cone coupling
\(\gamma = \big((x,y)\mapsto ((x,u(x,y)^{1/p}),(y,v(x,y)^{1/p}))\big)_\#\lambda.\)
Then $\mathsf P_p\gamma_1=\al$ and $\mathsf P_p\gamma_2=\be$ by the
weighted marginal constraints. Moreover
\(\int \De((x,r),(y,s))^p\d\gamma = \int H_{c(x,y)}(u(x,y),v(x,y))\d\lambda(x,y).\)
Taking the infimum gives $\CW\leq\HW$.

Conversely, let $\gamma$ be feasible for $\CW$. Choose arbitrary
representatives of cone points at the apex. This does not affect the
weighted constraints, and it does not change the cost because, by the
recession convention, $H_{c(x,y)}(0,s)$ and $H_{c(x,y)}(r,0)$ do not depend
on the arbitrary spatial representative attached to the zero radius. Let
$\lambda$ be the projection of $\gamma$ onto the spatial variables
$(x,y)$, and disintegrate
$\gamma=\gamma_{x,y}\d\lambda(x,y)$ with respect to this projection. Set
\(u(x,y)=\int r^p\d\gamma_{x,y}(r,s), \qquad v(x,y)=\int s^p\d\gamma_{x,y}(r,s).\)
The cone marginal constraints imply
$(\mathrm{p}_1)_\#(u\lambda)=\al$ and
$(\mathrm{p}_2)_\#(v\lambda)=\be$. Since $H_c$ is convex in its two
arguments, as a perspective of the convex reverse cost $L_c$, Jensen's
inequality gives, for $\lambda$-a.e. $(x,y)$, \(H_{c(x,y)}(u(x,y),v(x,y)) \leq \int H_{c(x,y)}(r^p,s^p)\d\gamma_{x,y}(r,s).\)
After integration and using
$\De((x,r),(y,s))^p=H_{c(x,y)}(r^p,s^p)$, this yields
$\HW\leq\CW$. Hence $\UW=\HW=\CW$.

Assume now that $\De$ is a distance. The nontrivial point in the triangle inequality is that two cone plans sharing the same weighted base projection need not have the same ordinary cone marginal. The radial homogeneity removes this obstruction. Indeed, common radial rescaling
\(((x,r),(y,s))\mapsto((x,r/\vartheta),(y,s/\vartheta)), \qquad \gamma\mapsto\vartheta^p\gamma,\)
preserves both weighted projections and the $p$-action. After adding mass at the cone apex, the normalization lemma for homogeneous lifts therefore lets one rescale two almost optimal plans for $(\alpha,\beta)$ and $(\beta,\zeta)$ so that their ordinary middle cone marginal is the same lift
\(\eta_\beta = M\delta_{\mathfrak o} + (x\mapsto(x,1))_\sharp\beta.\)
Here $M<+\infty$ is chosen large enough for both plans; additional mass at $(\mathfrak o,\mathfrak o)$ changes neither the weighted projections nor the action. The ordinary gluing lemma now produces a joint cone plan $(Z_0,Z_1,Z_2)$. Since $\De$ is a metric, Minkowski's inequality gives
\[
\left(\int\De(Z_0,Z_2)^p\right)^{1/p}
\leq
\left(\int\De(Z_0,Z_1)^p\right)^{1/p}
+
\left(\int\De(Z_1,Z_2)^p\right)^{1/p}.
\]
Taking infima proves the triangle inequality. This normalization-and-gluing argument is the metric engine of the homogeneous cone construction.

For completeness, normalize each non-apex pair by taking $\vartheta=(r^p+s^p)^{1/p}$. The transformed radii then satisfy $r^p+s^p=1$, while the transformed plan has total mass $\alpha(\Xx)+\beta(\Xx)$; apex mass can again be fixed without changing the problem. Thus minimizing plans may be restricted to a weakly compact family with bounded radii. Compactness of $\Xx$ and lower semicontinuity of the cone cost give an optimal plan. If $\CW(\alpha,\beta)=0$, such a plan is concentrated on the diagonal of the cone. There $x=y$ whenever the common radius is positive and $r=s$, so its two weighted base projections coincide and $\alpha=\beta$. Nonnegativity and symmetry are immediate, completing the metric proof.
\end{proof}
\paragraph{Entropic KL relaxation.}
\(\inf_{\pi\in\Mm_+(\Xx\times\Yy)} \int c\d\pi + \Dd_{\psi_1}(\pi_1|\al) + \Dd_{\psi_2}(\pi_2|\be) + \epsilon\Dd_\phi(\pi|\al\otimes\be).\)
\[
\sup_{f,g}
-
\Dd_{\psi_1}^*(-f|\al)
-
\Dd_{\psi_2}^*(-g|\be)
-
\epsilon\Dd_\phi^*\left(\frac{f\oplus g-c}{\epsilon}\middle|\al\otimes\be\right).
\]
For $\Dd_\phi=\KL$, the primal-dual relation is $\d\pi=e^{(f\oplus g-c)/\epsilon}\d\al\d\be$. If, in addition, $\Dd_{\psi_1}=\Dd_{\psi_2}=\tau\KL$, the dual specializes to

\[
\sup_{f,g}
-\tau\int(e^{-f/\tau}-1)\d\al
-\tau\int(e^{-g/\tau}-1)\d\be
-\epsilon\iint
\left(e^{(f\oplus g-c)/\epsilon}-1\right)
\d\al\d\be,
\]
\[
f \leftarrow \omega\, g^{\bar c,\epsilon},
\qquad
g \leftarrow \omega\, f^{c,\epsilon},
\qquad
\omega\eqdef\frac{\tau}{\tau+\epsilon},
\]
where the soft $c$-transforms are those of Definition~\ref{def-continuous-soft-c-transform}. Equivalently,

\begin{align*}
f(x)
&=
-
\frac{\tau\epsilon}{\tau+\epsilon}
\log\int_{\Yy}
\exp\left(\frac{g(y)-c(x,y)}{\epsilon}\right)\d\be(y),\\
g(y)
&=
-
\frac{\tau\epsilon}{\tau+\epsilon}
\log\int_{\Xx}
\exp\left(\frac{f(x)-c(x,y)}{\epsilon}\right)\d\al(x).
\end{align*}
\[
u_i\leftarrow \left(\frac{\a_i}{(K v)_i}\right)^\omega,
\qquad
v_j\leftarrow \left(\frac{\b_j}{(\transp{K}u)_j}\right)^\omega,
\qquad
\P=\diag(u)K\diag(v).
\]
\paragraph{Metric contraction of the damped updates.}
For balanced entropic OT, potentials are only defined up to the gauge $(f,g)\sim(f+\lambda,g-\lambda)$, so Hilbert's projective metric is natural. The KL marginal penalties break this invariance: under the usual finite-dual assumptions, the dual potentials are unique up to $\al$- and $\be$-null sets, and the natural distance is the ordinary $L^\infty$ distance on potentials.

\(d_T\big((u,v),(u',v')\big) \eqdef \max\{\norm{\log u-\log u'}_\infty,\norm{\log v-\log v'}_\infty\},\)
\begin{prop}[Linear contraction of unbalanced Sinkhorn]\label{prop-unbalanced-sinkhorn-contraction}
Assume that the damped soft-transform map sends bounded potentials to bounded potentials; for instance, this holds when \(c\) is bounded. Let
$\omega=\tau/(\tau+\epsilon)<1$, and define
\(d_\infty((f,g),(f',g')) \eqdef \max\{\norm{f-f'}_\infty,\norm{g-g'}_\infty\}.\)
The damped soft-transform map
\[
\mathcal T_{\rm GS}(f,g)
\eqdef
\big(\tilde f,\tilde g\big),
\qquad
\tilde f=\omega\,g^{\bar c,\epsilon},
\qquad
\tilde g=\omega\,\tilde f^{c,\epsilon},
\]
is a contraction: \(d_\infty\big(\mathcal T_{\rm GS}(f,g),\mathcal T_{\rm GS}(f',g')\big) \leq \omega\, d_\infty((f,g),(f',g')).\)
$(f^\star,g^\star)$ and the iterates satisfy
\(d_\infty((f^{(k)},g^{(k)}),(f^\star,g^\star)) \leq \omega^k d_\infty((f^{(0)},g^{(0)}),(f^\star,g^\star)).\)
Equivalently, for $u^{(k)}=e^{f^{(k)}/\epsilon}$ and
$v^{(k)}=e^{g^{(k)}/\epsilon}$, \(d_T((u^{(k)},v^{(k)}),(u^\star,v^\star)) \leq \omega^k d_T((u^{(0)},v^{(0)}),(u^\star,v^\star)).\)
\end{prop}
\begin{proof}
The soft $c$-transform is $1$-Lipschitz in $L^\infty$. Indeed, if
$\norm{g-g'}_\infty\leq \delta$, then $g-\delta\leq g'\leq g+\delta$.
Since $g\mapsto g^{\bar c,\epsilon}$ is order reversing and satisfies
$(g+\lambda)^{\bar c,\epsilon}=g^{\bar c,\epsilon}-\lambda$, this implies
\(g^{\bar c,\epsilon}-\delta \leq (g')^{\bar c,\epsilon} \leq g^{\bar c,\epsilon}+\delta.\)
Hence
$\norm{g^{\bar c,\epsilon}-(g')^{\bar c,\epsilon}}_\infty\leq
\norm{g-g'}_\infty$, and the same argument applies to
$f\mapsto f^{c,\epsilon}$.

Let $(\tilde f,\tilde g)=\mathcal T_{\rm GS}(f,g)$ and
$(\tilde f',\tilde g')=\mathcal T_{\rm GS}(f',g')$. With
$D=d_\infty((f,g),(f',g'))$, the previous Lipschitz estimate gives \(\norm{\tilde f-\tilde f'}_\infty \leq \omega\norm{g-g'}_\infty \leq \omega D,\)
and then
\(\norm{\tilde g-\tilde g'}_\infty \leq \omega\norm{\tilde f-\tilde f'}_\infty \leq \omega^2D \leq \omega D.\)
This proves the contraction. The Banach fixed-point theorem gives
existence, uniqueness and the displayed linear rate.
\end{proof}
\paragraph{Partial optimal transport.}
\(0\leq m\leq \min\{\al(\Xx),\be(\Yy)\},\)
\[
\operatorname{POT}_m(\al,\be)
\eqdef
\inf_{\substack{\pi\in\Mm_+(\Xx\times\Yy)\\
\pi_1\leq\al,\ \pi_2\leq\be\\
\pi(\Xx\times\Yy)=m}}
\int c\,\d\pi.
\]
Thus only a submeasure of $\al$ is transported onto a submeasure of $\be$; the remaining mass is left unmatched. The corresponding Lagrangian form is obtained by adding total-variation penalties.

\(\inf_{\pi\in\Mm_+(\Xx\times\Yy)} \int c\,\d\pi + \lambda\norm{\al-\pi_1}_{\TV} + \lambda\norm{\be-\pi_2}_{\TV}.\)
Assume $c\geq0$. Allowing transported marginals larger than the available marginals does not improve the value: excess transported mass can be trimmed, which decreases the transport cost and cannot increase the sum of the two total-variation penalties. Hence an optimal plan may be chosen with $\pi_1\leq\al$ and $\pi_2\leq\be$. If $m=\pi(\Xx\times\Yy)$, the penalty then reduces to

\(\lambda\big(\al(\Xx)-m\big)+\lambda\big(\be(\Yy)-m\big),\)
\begin{prop}[TV penalization selects fixed-mass partial OT]\label{prop-tv-partial-ot-lagrange}
Assume that $\Xx,\Yy$ are compact metric spaces, that $c$ is continuous and nonnegative, and set $A=\al(\Xx)$, $B=\be(\Yy)$, $M=\min(A,B)$. For $\lambda\geq0$, define the TV-penalized partial value
\[
\operatorname{POT}^{\TV}_\lambda(\al,\be)
\eqdef
\inf_{\substack{\pi\in\Mm_+(\Xx\times\Yy)\\ \pi_1\leq\al,\ \pi_2\leq\be}}
\left\{
\int c\,\d\pi
+
\lambda\big(A+B-2\pi(\Xx\times\Yy)\big)
\right\}.
\]
Then
\[
\operatorname{POT}^{\TV}_\lambda(\al,\be)
=
\lambda(A+B)
+
\inf_{0\leq m\leq M}
\left\{
\operatorname{POT}_m(\al,\be)-2\lambda m
\right\}.
\]
If $\pi_\lambda$ minimizes $\operatorname{POT}^{\TV}_\lambda$ and $m_\lambda=\pi_\lambda(\Xx\times\Yy)$, then $m_\lambda$ minimizes $m\mapsto\operatorname{POT}_m(\al,\be)-2\lambda m$ and $\pi_\lambda$ minimizes $\operatorname{POT}_{m_\lambda}(\al,\be)$. Conversely, if $2\lambda$ is a subgradient at $m$ of the convex function $m\mapsto\operatorname{POT}_m(\al,\be)$ extended by $+\infty$ outside $[0,M]$, then every minimizer of $\operatorname{POT}_m(\al,\be)$ is a minimizer of $\operatorname{POT}^{\TV}_\lambda(\al,\be)$. In particular, every prescribed mass $m\in[0,M]$ is selected by at least one multiplier $\lambda\geq0$, possibly together with a whole interval of masses when the value function has a flat exposed face.
\end{prop}
\begin{proof}
Write $V(m)=\operatorname{POT}_m(\al,\be)$. Compactness and continuity give an optimizer for every $m\in[0,M]$: the admissible subcouplings form a weakly compact set and the cost is continuous. The function $V$ is convex on $[0,M]$: if $\pi_0,\pi_1$ are admissible competitors of masses $m_0,m_1$, then $(1-t)\pi_0+t\pi_1$ is admissible for mass $(1-t)m_0+tm_1$ and has the corresponding convex-combination cost. It is also nondecreasing because any competitor of mass $m_1>m_0$ can be rescaled by $m_0/m_1$ to give a competitor of mass $m_0$ with no larger cost. Finally, $V$ is Lipschitz. If $0\leq m_0<m_1\leq M$ and $\pi_0$ is optimal at mass $m_0$, both residual measures $\alpha-(\pi_0)_1$ and $\beta-(\pi_0)_2$ have mass at least $m_1-m_0$. Choose submeasures of each residual with this common mass and couple them. Adding the resulting plan to $\pi_0$ gives an admissible plan of mass $m_1$ at an additional cost bounded by $\|c\|_\infty(m_1-m_0)$.

For any admissible $\pi$ in the definition of $\operatorname{POT}^{\TV}_\lambda$, let $m=\pi(\Xx\times\Yy)$. The plan is admissible for $V(m)$, hence its objective is at least $V(m)+\lambda(A+B-2m)$. Taking the infimum over $\pi$ gives the lower bound by the displayed right-hand side. Conversely, for each $m$ and each minimizing sequence for $V(m)$, the same plans are admissible in $\operatorname{POT}^{\TV}_\lambda$ and have objective tending to $V(m)+\lambda(A+B-2m)$. Taking the infimum over $m$ gives the reverse inequality.

The optimizer statement follows from the same inequalities. If $\pi_\lambda$ minimizes the penalized problem with mass $m_\lambda$, then equality in the lower bound forces both $V(m_\lambda)-2\lambda m_\lambda$ to be minimal over $[0,M]$ and $\pi_\lambda$ to be optimal at fixed mass $m_\lambda$. Conversely, $2\lambda\in\partial V(m)$ means
\(V(m')-2\lambda m'\geq V(m)-2\lambda m \quad\text{for all }m'\in[0,M].\)
Thus any minimizer of $V(m)$ attains the penalized infimum. Because the extended function is proper, convex, lower semicontinuous and Lipschitz on its effective interval, its subdifferential is nonempty at every $m\in[0,M]$ after adding the normal cone of the interval. Monotonicity of $V$ ensures that a subgradient can be chosen nonnegative, hence equal to $2\lambda$ for some $\lambda\geq0$.
\end{proof}
\begin{prop}[Fixed total mass reduces partial OT to balanced OT]\label{prop-partial-ot-metric-slice}
Let $(\Xx,d)$ be a metric space, let $p\geq1$, set $c=d^p$, and fix $m>0$. On the class of nonnegative measures with total mass exactly $m$ and finite $p$th moment,
\(\mathcal M_m^p(\Xx)\eqdef \enscond{\al\in\Mm_+(\Xx)}{\al(\Xx)=m,\ \int d(x_0,x)^p\d\al(x)<+\infty}\)
for one, hence any, base point $x_0$, the quantity
\(d_m(\al,\be)\eqdef \operatorname{POT}_m(\al,\be)^{1/p}\)
is a distance. More precisely, if $\bar\al=\al/m$ and $\bar\be=\be/m$, then
\(d_m(\al,\be)=m^{1/p}\Wass_p(\bar\al,\bar\be).\)
If instead $\operatorname{POT}_m$ is considered on a larger class of measures with mass at least $m$, then it is not a distance in general for $m>0$: different measures can share a common submeasure of mass $m$, giving zero partial cost.
\end{prop}
\begin{proof}
If $\al(\Xx)=\be(\Xx)=m$ and $\pi$ is admissible in $\operatorname{POT}_m(\al,\be)$, then $\pi_1\leq\al$ and both measures have mass $m$, hence $\pi_1=\al$. Similarly $\pi_2=\be$. Thus the admissible plans are exactly the ordinary couplings between $\al$ and $\be$. Dividing such a coupling by $m$ gives a probability coupling between $\bar\al$ and $\bar\be$, and conversely multiplying any coupling of $\bar\al,\bar\be$ by $m$ gives a coupling of $\al,\be$. Therefore
\(\operatorname{POT}_m(\al,\be) = m\Wass_p(\bar\al,\bar\be)^p.\)
The metric properties of $d_m$ follow from those of $\Wass_p$ in Proposition~\ref{prop-metric-measure}.

For the last claim, take three distinct points $x,y,z$ in a metric space and a number $r>0$. The measures $\al=m\de_x+r\de_y$ and $\be=m\de_x+r\de_z$ are different but the partial plan $m\de_{(x,x)}$ is admissible and has zero cost. Hence $\operatorname{POT}_m(\al,\be)=0$ although $\al\neq\be$.
\end{proof}
Thus the TV-penalized formulation is the Lagrangian envelope of constrained partial OT: increasing $\lambda$ penalizes unmatched mass more strongly and therefore selects larger transported masses, while small $\lambda$ makes deletion and creation cheaper. As the prescribed transported mass decreases, the plan keeps only the least costly facing portions of the two marginals.

\subsection{Sliced Wasserstein Distances}
\label{sec-sliced-wasserstein}
\paragraph{One-dimensional projections.}
\paragraph{Spherical averaging.}
\begin{defn}[Sliced Wasserstein distance]\label{def-sliced-wasserstein}
Let $p\geq1$, let $\alpha,\beta\in\Pp_p(\RR^d)$, and let $\sigma$ be the uniform probability measure on the sphere $\Sphere^{d-1}$. The sliced $p$-Wasserstein distance is
\eql{\label{eq-sliced-wasserstein}
\SW_p(\al,\be)^p
\eqdef
\int_{\Sphere^{d-1}}
\Wass_p\big((P_\theta)_\sharp\al,(P_\theta)_\sharp\be\big)^p
\d\sigma(\theta).
}
\end{defn}
\begin{prop}[Metric properties of sliced Wasserstein]\label{prop-sliced-wasserstein-metric}
For $p\geq1$, $\SW_p$ is a distance on $\Pp_p(\RR^d)$. Moreover, $\SW_p$ metrizes weak convergence together with convergence of the $p$th moment. Finally, \(\SW_p(\alpha,\beta)\leq \Wass_p(\alpha,\beta),\)
and, for $p=2$ with the uniform probability measure on the sphere, \(\SW_2(\alpha,\beta)^2\leq \frac1d\Wass_2(\alpha,\beta)^2.\)
Conversely, if $d\geq2$ and $\alpha,\beta\in\Pp_1(\RR^d)$ are supported in
\(B_R\eqdef\enscond{x\in\RR^d}{\norm{x}\leq R}\), then there is a constant
$C_d$ depending only on $d$ such that
\(\Wass_1(\alpha,\beta) \leq C_d R^{\frac{d}{d+1}}\SW_1(\alpha,\beta)^{\frac{1}{d+1}}.\)
Moreover, the sharp dimension-dependent comparison for $p=1$ has the form
\(\Wass_1(\alpha,\beta) \leq c_d R^{\frac{d-1}{d}}\SW_1(\alpha,\beta)^{\frac1d}.\)
\end{prop}
\begin{proof}
Non-negativity and symmetry follow from the one-dimensional Wasserstein distance. For the triangle inequality, apply the triangle inequality of $\Wass_p$ for each direction $\theta$ and then Minkowski's inequality in $L^p(\Sphere^{d-1})$.

If $\SW_p(\alpha,\beta)=0$, then $(P_\theta)_\sharp\alpha=(P_\theta)_\sharp\beta$ for almost every direction. By continuity of characteristic functions this holds for all directions, and the Cram\'er--Wold theorem implies $\alpha=\beta$. This proves separation.

The bound $\SW_p\leq\Wass_p$ follows because $P_\theta$ is $1$-Lipschitz, so
\(\Wass_p((P_\theta)_\sharp\alpha,(P_\theta)_\sharp\beta) \leq \Wass_p(\alpha,\beta)\)
for every $\theta$. For $p=2$, using any coupling $\pi$ between $\alpha$ and $\beta$, \(\int_{\Sphere^{d-1}}\int |\dotp{x-y}{\theta}|^2\d\pi(x,y)\d\sigma(\theta) = \frac1d\int\norm{x-y}^2\d\pi(x,y).\)
Optimizing over $\pi$ gives the sharper inequality.

To prove the topological statement, first suppose that $\Wass_p(\alpha_n,\alpha)\to0$. The inequality $\SW_p\leq\Wass_p$ immediately gives $\SW_p(\alpha_n,\alpha)\to0$. Conversely, suppose that $\SW_p(\alpha_n,\alpha)\to0$. The spherical identity
\[
\int_{\Sphere^{d-1}}\int_{\RR^d}|\langle\theta,x\rangle|^p\d\alpha_n(x)\d\sigma(\theta)
=
c_{d,p}\int_{\RR^d}\|x\|^p\d\alpha_n(x),
\qquad
c_{d,p}\eqdef\int_{\Sphere^{d-1}}|\theta_1|^p\d\sigma(\theta)>0,
\]
and the triangle inequality in $L^p(\Sphere^{d-1}\times(0,1))$ show that the $p$th moments of $(\alpha_n)$ are uniformly bounded. Hence the sequence is tight. From every subsequence one can extract a further subsequence for which the projected $\Wass_p$ distances converge to zero for $\sigma$-a.e. $\theta$. Every weak limit therefore has the same one-dimensional projections as $\alpha$, so Cram\'er--Wold identifies it with $\alpha$. Thus $\alpha_n\rightharpoonup\alpha$. Finally, let $Q_{n,\theta}$ and $Q_\theta$ be the quantile functions of the projected measures. The one-dimensional quantile formula gives
\(\SW_p(\alpha_n,\alpha) = \|Q_{n,\theta}-Q_\theta\|_{L^p(\Sphere^{d-1}\times(0,1))}.\)
Hence the corresponding $L^p$ norms converge. By the spherical identity, their $p$th powers are respectively $c_{d,p}\int\|x\|^p\d\alpha_n(x)$ and $c_{d,p}\int\|x\|^p\d\alpha(x)$. Thus the $p$th moments converge. Weak convergence plus convergence of these moments is equivalent to convergence in $\Wass_p$.

The reverse estimates on compact sets are deeper and use the Radon structure of slices. Bonnotte's Lemma~5.1.4 gives the historical exponent $1/(d+1)$ by mollifying a Kantorovich--Rubinstein test function, representing the regularized test through one-dimensional projections, bounding the resulting action by $\SW_1$, and optimizing the smoothing scale. The scaling in $R$ follows by applying the unit-ball estimate to $(x\mapsto x/R)_\sharp\alpha$ and $(x\mapsto x/R)_\sharp\beta$. The exponent $1/d$ is the sharp comparison theorem of Carlier, Figalli, M\'erigot and Wang; their examples show that it cannot be improved in general.
\end{proof}
\begin{prop}[First-order comparison along Brenier perturbations]\label{prop-sliced-first-order-tangent}
Let $\alpha\in\Pp_2(\RR^d)$ be absolutely continuous, and let $\varphi\in C^2(\RR^d)$ be such that $\nabla\varphi\in L^2(\alpha;\RR^d)$ and $\nabla^2\varphi$ is bounded. For $\abs{t}$ small enough, set
\(T_t(x)\eqdef x+t\nabla\varphi(x), \qquad \alpha_t\eqdef(T_t)_\sharp\alpha.\)
Then $T_t$ is the quadratic Brenier map from $\alpha$ to $\alpha_t$ for $\abs{t}$ small enough, and
\(\Wass_2(\alpha,\alpha_t)^2 = t^2\int_{\RR^d}\norm{\nabla\varphi(x)}^2\d\alpha(x).\)
Moreover, \(\SW_2(\alpha,\alpha_t)^2 \leq \frac{t^2}{d}\int_{\RR^d}\norm{\nabla\varphi(x)}^2\d\alpha(x).\)
Thus the sliced infinitesimal speed is at most $d^{-1/2}$ times the Wasserstein infinitesimal speed along smooth Brenier perturbations.
\end{prop}
\begin{proof}
The map $T_t$ is the gradient of
\(\psi_t(x)\eqdef \frac12\norm{x}^2+t\varphi(x).\)
Since $\nabla^2\varphi$ is bounded, $\nabla^2\psi_t=\Id+t\nabla^2\varphi$ is positive semidefinite for $\abs{t}$ small enough. Hence $\psi_t$ is convex and $T_t=\nabla\psi_t$ is a Brenier map. Its graph is cyclically monotone, so the Brenier optimality criterion gives the optimality of $(\Id,T_t)_\sharp\alpha$, and
\(\Wass_2(\alpha,\alpha_t)^2 = \int\norm{x-T_t(x)}^2\d\alpha(x) = t^2\int\norm{\nabla\varphi(x)}^2\d\alpha(x).\)
For a fixed direction $\theta$, the coupling
\(\big(P_\theta,P_\theta\circ T_t\big)_\sharp\alpha\)
is admissible between $(P_\theta)_\sharp\alpha$ and $(P_\theta)_\sharp\alpha_t$. Therefore
\(\Wass_2\big((P_\theta)_\sharp\alpha,(P_\theta)_\sharp\alpha_t\big)^2 \leq t^2\int\abs{\dotp{\theta}{\nabla\varphi(x)}}^2\d\alpha(x).\)
Integrating over the sphere and using
\(\int_{\Sphere^{d-1}}\abs{\dotp{\theta}{w}}^2\d\sigma(\theta) = \frac1d\norm{w}^2\)
gives the sliced estimate.
\end{proof}
\begin{prop}[Translations separate sliced length from $\Wass_2$]\label{prop-sliced-length-translation}
Let $\tau_h(x)=x+h$ for $h\in\RR^d$. For any $\alpha\in\Pp_2(\RR^d)$, \(\SW_2^{\mathrm{len}}\big(\alpha,(\tau_h)_\sharp\alpha\big) = \SW_2\big(\alpha,(\tau_h)_\sharp\alpha\big) = \frac{\norm{h}}{\sqrt d},\)
whereas
\(\Wass_2\big(\alpha,(\tau_h)_\sharp\alpha\big)=\norm{h}.\)
\end{prop}
\begin{proof}
For each direction $\theta$, the projected measure $(P_\theta)_\sharp(\tau_h)_\sharp\alpha$ is the translate of $(P_\theta)_\sharp\alpha$ by the scalar $\dotp{\theta}{h}$. Hence
\(\Wass_2\big((P_\theta)_\sharp\alpha,(P_\theta)_\sharp(\tau_h)_\sharp\alpha\big)^2 = \abs{\dotp{\theta}{h}}^2.\)
After integration on the sphere, \(\SW_2\big(\alpha,(\tau_h)_\sharp\alpha\big)^2 = \int_{\Sphere^{d-1}}\abs{\dotp{\theta}{h}}^2\d\sigma(\theta) = \frac{\norm{h}^2}{d}.\)
For the translation curve $\alpha_s=(\tau_{sh})_\sharp\alpha$, the same computation gives \(\SW_2(\alpha_s,\alpha_t)=\abs{s-t}\frac{\norm{h}}{\sqrt d}.\)
Every partition therefore has $\SW_2$-length exactly $\norm{h}/\sqrt d$, which gives the upper bound on $\SW_2^{\mathrm{len}}$. The opposite inequality follows from the general fact that any metric distance is bounded by the length of every connecting curve.

Finally, the transport map $\tau_h=\nabla(\frac12\norm{x}^2+\dotp{h}{x})$ is the gradient of a convex function. By the Brenier cyclic-monotonicity criterion, the translation coupling is optimal for the quadratic cost, and hence the $\Wass_2$ distance is $\norm{h}$.
\end{proof}
\begin{defn}[Max-sliced Wasserstein]\label{def-sliced-variants}
The max-sliced distance replaces the average over directions by the most discriminating one: \(\MaxSW_p(\alpha,\beta) \eqdef \sup_{\theta\in\Sphere^{d-1}} \Wass_p((P_\theta)_\sharp\alpha,(P_\theta)_\sharp\beta).\)
\end{defn}
\paragraph{Subspace-sliced variants.}
\begin{defn}[Subspace-sliced Wasserstein]\label{def-subspace-sliced-wasserstein}
Fix $1\leq k\leq d$. Subspace-sliced variants replace one-dimensional lines by $k$-dimensional orthogonal projections. If $U\in\RR^{d\times k}$ satisfies $\transp{U}U=\Id_k$, then
\(\SW_{p,k}(\alpha,\beta)^p \eqdef \int \Wass_p((\transp{U})_\sharp\alpha,(\transp{U})_\sharp\beta)^p\d U,\)
where $\d U$ denotes the normalized invariant measure on the Stiefel manifold
$\mathrm{St}(d,k)=\{U\in\RR^{d\times k}\;:\;\transp{U}U=\Id_k\}$, and
\(\MaxSW_{p,k}(\alpha,\beta) \eqdef \sup_{\transp{U}U=\Id_k} \Wass_p((\transp{U})_\sharp\alpha,(\transp{U})_\sharp\beta).\)
The case $k=1$ recovers ordinary sliced and max-sliced Wasserstein, while $k=d$ recovers the original Wasserstein distance.
\end{defn}
\begin{prop}[Basic bounds for sliced variants]\label{prop-sliced-variant-bounds}
Let $p\geq1$ and let $\alpha,\beta\in\Pp_p(\RR^d)$. With normalized spherical and Stiefel measures, \(\SW_p(\alpha,\beta) \leq \MaxSW_p(\alpha,\beta) \leq \Wass_p(\alpha,\beta).\)
For $k$-dimensional subspace projections, \(\SW_{p,k}(\alpha,\beta) \leq \MaxSW_{p,k}(\alpha,\beta) \leq \Wass_p(\alpha,\beta).\)
For $p=2$, the averaged subspace distance satisfies the sharper estimate
\(\SW_{2,k}(\alpha,\beta)^2 \leq \frac{k}{d}\Wass_2(\alpha,\beta)^2.\)
\end{prop}
\begin{proof}
The first inequality in each line follows because an $L^p$ average over a probability space is bounded by the corresponding supremum. The second inequality follows because orthogonal projections are $1$-Lipschitz: pushing any admissible coupling between $\alpha$ and $\beta$ through a projection gives an admissible coupling for the projected measures with no larger transport cost. Optimizing over couplings and then averaging or maximizing over the projection gives the result. For $p=2$, rotational invariance of the Stiefel measure gives
\(\int_{\mathrm{St}(d,k)}\|U^\top z\|^2\d U = \frac{k}{d}\|z\|^2.\)
Apply this identity to $z=x-y$ under any coupling and then minimize over couplings.
\end{proof}
\paragraph{Min-sliced lifted transport plans.}
\(\pi_\theta = \frac1n\sum_{i=1}^n\delta_{(x_i,y_{\sigma_\theta(i)})}\)
\(\operatorname{MSWGG}_2(\alpha,\beta)^2 \eqdef \min_{\theta\in\Sphere^{d-1}} \int\norm{x-y}^2\d\pi_\theta(x,y).\)
\(\Wass_2^2(\alpha,\beta) \leq \int\norm{x-y}^2\d\pi_\theta(x,y), \qquad \Wass_2^2(\alpha,\beta) \leq \operatorname{MSWGG}_2(\alpha,\beta)^2.\)
\subsection{Quotient Wasserstein and Wasserstein--Procrustes}
\label{sec-quotient-wasserstein-procrustes}
\paragraph{Quotient distances.}
\begin{defn}[Quotient Wasserstein distance]\label{def-quotient-wasserstein}
Let a group $\mathcal G$ act on a metric space $(\Xx,d)$ by isometries, and assume that the action preserves finite $p$th moments. Write $g_\sharp\alpha$ for the push-forward of a measure $\alpha$ by $g\in\mathcal G$, and $[\alpha]\eqdef\{g_\sharp\alpha\;:\;g\in\mathcal G\}$ for its orbit. The quotient $p$-Wasserstein distance between the orbits of $\alpha,\beta\in\Pp_p(\Xx)$ is
\[
\Wass_{p,\mathcal G}([\alpha],[\beta])
\eqdef
\inf_{g,h\in\mathcal G}\Wass_p(g_\sharp\alpha,h_\sharp\beta)
=
\inf_{g\in\mathcal G}\Wass_p(\alpha,g_\sharp\beta).
\]
\end{defn}
\begin{prop}[Metric property on the quotient]\label{prop-quotient-wasserstein-metric}
If $\mathcal G$ acts by isometries and preserves $\Pp_p(\Xx)$, then $\Wass_{p,\mathcal G}$ is a pseudo-distance on $\Pp_p(\Xx)$. If the infimum is attained between every pair of orbits, then it is a distance on the orbit space $\Pp_p(\Xx)/\mathcal G$. This applies in particular to continuous actions of compact groups.
\end{prop}
\begin{proof}
Isometry of the action gives $\Wass_p(g_\sharp\alpha,g_\sharp\beta)=\Wass_p(\alpha,\beta)$, which proves the second formula in Definition~\ref{def-quotient-wasserstein}. Non-negativity and symmetry are inherited from $\Wass_p$. For the triangle inequality, choose $g_1,g_2\in\mathcal G$. Then
\[
\Wass_p(\alpha,(g_1)_\sharp\beta)
+
\Wass_p(\beta,(g_2)_\sharp\gamma)
=
\Wass_p(\alpha,(g_1)_\sharp\beta)
+
\Wass_p((g_1)_\sharp\beta,(g_1g_2)_\sharp\gamma).
\]
Applying the triangle inequality for $\Wass_p$ and then taking infima over $g_1,g_2$ gives the result. If the infimum is attained and the quotient distance is zero, there exist $g,h\in\mathcal G$ such that $\Wass_p(g_\sharp\alpha,h_\sharp\beta)=0$, hence $g_\sharp\alpha=h_\sharp\beta$ and therefore $[\alpha]=[\beta]$. For a continuous action of a compact group, the objective is continuous on a compact parameter set, so attainment follows. Without attainment, the construction is naturally a metric only after identifying distinct orbits at zero orbit distance.
\end{proof}
\paragraph{Rigid motions and Wasserstein--Procrustes.}
\(\inf_{R\in\mathrm O(d),\,t\in\RR^d,\,\pi\in\Couplings(\alpha,\beta)} \int\norm{Rx+t-y}^2\d\pi(x,y).\)
This is an extrinsic counterpart of the Gromov--Wasserstein viewpoint developed in Section~\ref{sec-gromov-wasserstein}. Procrustes alignment searches only over ambient rigid motions, whereas GW is invariant under all measure-preserving isometries of the metric-measure spaces.

\(\GW((\RR^d,\norm{\cdot},\alpha),(\RR^d,\norm{\cdot},\beta)) = \GW((\RR^d,\norm{\cdot},g_\sharp\alpha),(\RR^d,\norm{\cdot},h_\sharp\beta)).\)
\(\GW((\RR^d,\norm{\cdot},g_\sharp\alpha),(\RR^d,\norm{\cdot},h_\sharp\beta)) \leq 2\,\Wass_p(g_\sharp\alpha,h_\sharp\beta).\)
\begin{equation}\label{eq-gw-procrustes-upper}
\GW((\RR^d,\norm{\cdot},\alpha),(\RR^d,\norm{\cdot},\beta))
\leq
2\,\Wass_{p,\mathrm E(d)}([\alpha],[\beta]),
\end{equation}
for the same exponent $p$ and $\De(u,v)=|u-v|$. Thus a good Wasserstein--Procrustes registration forces small GW distortion.

\begin{equation}\label{eq-wasserstein-procrustes-coupling-update}
\P^{(k)}
\in
\argmin_{\P\in\CouplingsD(\a,\b)}
\sum_{i,j}\P_{ij}\norm{R^{(k)}x_i+t^{(k)}-y_j}^2.
\end{equation}
\begin{equation}\label{eq-wasserstein-procrustes-rigid-update}
(R^{(k+1)},t^{(k+1)})
\in
\argmin_{R\in\mathrm O(d),\,t\in\RR^d}
\sum_{i,j}\P^{(k)}_{ij}\norm{Rx_i+t-y_j}^2,
\end{equation}
\begin{prop}[Rigid update for a fixed coupling]\label{prop-wasserstein-procrustes-rigid-update}
Let $\alpha=\sum_i a_i\delta_{x_i}$ and $\beta=\sum_j b_j\delta_{y_j}$, and fix $\P\in\CouplingsD(\a,\b)$. Define
\(\bar x=\sum_i a_i x_i,\qquad \bar y=\sum_j b_j y_j,\qquad M_\P=\sum_{i,j} \P_{ij}(y_j-\bar y)(x_i-\bar x)^\top.\)
If $M_\P=U\Sigma V^\top$ is a singular value decomposition, then the minimizers over the full orthogonal group of
\(\min_{R\in\mathrm O(d),\,t\in\RR^d} \sum_{i,j}\P_{ij}\norm{Rx_i+t-y_j}^2\)
are given by $R^\star=UV^\top$ and $t^\star=\bar y-R^\star\bar x$, up to the usual non-uniqueness when $M_\P$ is rank deficient. If one imposes $R\in\mathrm{SO}(d)$, set
\(D=\diag(1,\ldots,1,\det(UV^\top)), \qquad R^\star=UDV^\top, \qquad t^\star=\bar y-R^\star\bar x.\)
\end{prop}
\begin{proof}
For fixed $R$, differentiating with respect to $t$ gives $t=\bar y-R\bar x$. Substituting this value centers the variables and leaves
\(\sum_{i,j}\P_{ij}\norm{R(x_i-\bar x)-(y_j-\bar y)}^2.\)
The two quadratic terms independent of $R$ are fixed, so the problem is equivalent to maximizing
\(\sum_{i,j}\P_{ij}\dotp{R(x_i-\bar x)}{y_j-\bar y} = \tr(R^\top M_\P).\)
Von Neumann's trace inequality gives $\tr(R^\top U\Sigma V^\top)\leq\tr(\Sigma)$ for $R\in\mathrm O(d)$, with equality for $R=UV^\top$. Under the constraint $R\in\mathrm{SO}(d)$, the same argument applies with the additional determinant constraint on $U^\top R V$, which gives the displayed diagonal correction.
\end{proof}
\subsection{Vector Quantiles and Linear Optimal Transport}
\label{sec-linear-ot}
\label{sec-vector-quantiles-linearized-transport}
\paragraph{Vector quantiles.}
Assume that $\rho$ is absolutely continuous. For a target law $\al$ with finite second moment, its vector quantile relative to $\rho$ is the Brenier map

\(\T_\al=\nabla\phi_\al, \qquad (\T_\al)_\sharp\rho=\al,\)
\(\min_{\T_\sharp\rho=\al} \int \norm{x-\T(x)}^2\d\rho(x).\)
\paragraph{Linearized Wasserstein coordinates.}
\begin{equation}\label{eq-lot-embedding}
\alpha \mapsto T_\alpha-\Id\in L^2(\rho;\RR^d),
\qquad
\LOT_\rho(\alpha,\beta)=\norm{T_\alpha-T_\beta}_{L^2(\rho)}.
\end{equation}
If one of the two targets equals the reference, the linearized distance is exact: for instance, $\LOT_\rho(\rho,\alpha)=\norm{T_\alpha-\Id}_{L^2(\rho)}=\Wass_2(\rho,\alpha)$. Uniqueness of the Brenier maps shows that $\LOT_\rho$ is itself a distance on the target measures for which these maps are defined.

\(\bar T=\sum_s\lambda_s T_{\alpha_s}, \qquad \bar\alpha_{\LOT}=\bar T_\sharp\rho.\)
This is exact in one dimension, where quantile functions linearize $\Wass_2$, and it is especially useful when many barycenters with changing weights must be evaluated quickly.

\paragraph{Principal components in linear OT coordinates.}
\(z_i \eqdef T_{\alpha_i}-\Id, \qquad \bar z \eqdef \frac1N\sum_{i=1}^N z_i.\)
\(\mathcal C_{\LOT} h \eqdef \frac1N\sum_{i=1}^N \dotp{z_i-\bar z}{h}_{L^2(\rho)} (z_i-\bar z),\)
\(e_k=\frac{1}{\sqrt{N\lambda_k}}\sum_{i=1}^N v_i^{(k)}(z_i-\bar z)\)
\(a_{i,k} \eqdef \dotp{z_i-\bar z}{e_k}_{L^2(\rho)}.\)
\(T_a(x) \eqdef x+\bar z(x)+\sum_{k=1}^m a_k e_k(x), \qquad \alpha_a \eqdef (T_a)_\sharp\rho.\)
\begin{prop}[Quantitative stability of linear OT]\label{prop-linear-ot-stability}
Let $X,Y\subset\RR^d$ be compact convex sets, assume that $X$ has positive Lebesgue measure, and let $\rho$ be normalized Lebesgue measure on $X$. Then there exists a constant $C=C(d,X,Y)$ such that, for all $\alpha,\beta\in\Pp(Y)$,
\(\Wass_2(\alpha,\beta)\leq \LOT_\rho(\alpha,\beta) \quad\text{and}\quad \LOT_\rho(\alpha,\beta) \leq C\Wass_1(\alpha,\beta)^{2/15}.\)
\end{prop}
\begin{proof}
The first inequality is immediate: $(T_\alpha,T_\beta)_\sharp\rho$ is a feasible coupling between $\alpha$ and $\beta$. The second inequality is the quantitative stability theorem of M\'erigot, Delalande and Chazal. Its proof first controls the difference of the two Kantorovich potentials in terms of $\Wass_1(\alpha,\beta)$, then combines convexity and interpolation estimates to control the gradients in $L^2(\rho)$; the resulting dimension-independent H\"older exponent is $2/15$.
In one dimension, quantile functions make the embedding isometric and the exponent improves to one. In several dimensions, the theorem is global but only H\"older; it should not be read as a global Lipschitz estimate in $\Wass_2$.
\end{proof}
\subsection{Spectral and Robust Wasserstein Distances}
\label{sec-spectral-subspace-wasserstein}
\begin{defn}[Monotone spectral gauge]\label{def-monotone-spectral-gauge}
A monotone spectral gauge on positive semidefinite matrices is a convex, positively $1$-homogeneous map $\gamma:\mathbb S_+^d\to\RR_+$ such that $\gamma(M)=0$ only for $M=0$, $\gamma(QM\transp{Q})=\gamma(M)$ for every orthogonal matrix $Q$, and
\(0\preceq M\preceq N \quad\Longrightarrow\quad \gamma(M)\leq\gamma(N).\)
\end{defn}
\begin{defn}[Spectral Wasserstein distance]\label{def-spectral-wasserstein}
Let $\gamma$ be a monotone spectral gauge. For a coupling $\pi\in\Couplings(\alpha,\beta)$, define its displacement covariance
\(M_\pi\eqdef\int (x-y)(x-y)^\top\d\pi(x,y).\)
The spectral Wasserstein distance associated with $\gamma$ is
\eql{\label{eq-spectral-wasserstein}
\Wass_\gamma(\alpha,\beta)^2
\eqdef
\inf_{\pi\in\Couplings(\alpha,\beta)}\gamma(M_\pi).
}
\end{defn}
The special case $\gamma(M)=\tr(M)$ gives the usual quadratic Wasserstein distance $\Wass_2$. The spectral gauge $\gamma(M)=\lambda_{\max}(M)$ instead measures the worst transported variance direction.

\eql{\label{eq-quadratic-projected-cost}
\Wass_{2,A}(\alpha,\beta)^2
\eqdef
\inf_{\pi\in\Couplings(\alpha,\beta)}
\int (x-y)^\top A(x-y)\d\pi(x,y)
=
\Wass_2((A^{1/2})_\sharp\alpha,(A^{1/2})_\sharp\beta)^2.
}
\eql{\label{eq-spectral-polar-set}
\mathcal B_\gamma
\eqdef
\enscond{A\succeq0}{\tr(AM)\leq\gamma(M)\quad\text{for all }M\succeq0},
}
\begin{prop}[Robust representation and metric equivalence]\label{prop-spectral-wasserstein-robust}
Assume, for simplicity, that the measures are compactly supported and that $\gamma$ is closed and finite on the positive semidefinite cone. Then
\(\Wass_\gamma(\alpha,\beta)^2 = \sup_{A\in\mathcal B_\gamma} \Wass_{2,A}(\alpha,\beta)^2.\)
If there exist constants $0<a\leq b<+\infty$ such that $a\Id\in\mathcal B_\gamma$ and $\mathcal B_\gamma\subset\enscond{A}{0\preceq A\preceq b\Id}$, equivalently
\(a\tr(M)\leq\gamma(M)\leq b\tr(M) \qquad (M\succeq0),\)
then
\(\sqrt a\,\Wass_2(\alpha,\beta) \leq \Wass_\gamma(\alpha,\beta) \leq \sqrt b\,\Wass_2(\alpha,\beta).\)
In particular, $\Wass_\gamma$ is a distance. When $\gamma$ is the restriction of a norm to the positive semidefinite cone, these bounds hold automatically in finite dimension, so $\Wass_\gamma$ is equivalent to $\Wass_2$ on measures with finite second moments.
\end{prop}
\begin{proof}
Using the polar representation of $\gamma$, \(\Wass_\gamma(\alpha,\beta)^2 = \inf_{\pi\in\Couplings(\alpha,\beta)} \sup_{A\in\mathcal B_\gamma}\tr(AM_\pi).\)
The coupling set is convex and compact for weak convergence under compact support. Since $\gamma$ is a finite gauge on a finite-dimensional cone and vanishes only at the origin, it is equivalent to the trace norm on the slice $\tr(M)=1$, so $\mathcal B_\gamma$ is convex and compact. The map $(\pi,A)\mapsto\tr(AM_\pi)$ is affine in each variable and continuous. Sion's minimax theorem gives
\[
\inf_\pi\sup_{A\in\mathcal B_\gamma}\tr(AM_\pi)
=
\sup_{A\in\mathcal B_\gamma}\inf_\pi\tr(AM_\pi)
=
\sup_{A\in\mathcal B_\gamma}\Wass_{2,A}(\alpha,\beta)^2.
\]
For fixed $A\succeq0$, $\Wass_{2,A}$ is the Wasserstein pseudodistance associated with the seminorm $x\mapsto\norm{A^{1/2}x}$. Since all terms are nonnegative, the robust identity also gives
\(\Wass_\gamma(\alpha,\beta) = \sup_{A\in\mathcal B_\gamma}\Wass_{2,A}(\alpha,\beta).\)
A supremum of pseudodistances is symmetric and satisfies the triangle inequality.

If $a\Id\in\mathcal B_\gamma$ and $A\preceq b\Id$ for all $A\in\mathcal B_\gamma$, then
\(a\Wass_2(\alpha,\beta)^2 = \Wass_{2,a\Id}(\alpha,\beta)^2 \leq \Wass_\gamma(\alpha,\beta)^2 \leq b\Wass_2(\alpha,\beta)^2,\)
which proves definiteness and equivalence with $\Wass_2$. The equivalence between these operator bounds and $a\tr(M)\leq\gamma(M)\leq b\tr(M)$ follows directly from the polar formula. In finite dimension, any norm restricted to the positive semidefinite cone is equivalent to the trace norm on that cone. The finite-second-moment case follows by truncation when these norm-equivalence bounds hold.
\end{proof}
\begin{defn}[Subspace robust Wasserstein]\label{def-subspace-robust-wasserstein}
For $1\leq k\leq d$, the Paty--Cuturi subspace robust Wasserstein distance is
\[
\SRW_{2,k}(\alpha,\beta)
\eqdef
\sup_{\transp{U}U=\Id_k}
\Wass_2((\transp{U})_\sharp\alpha,(\transp{U})_\sharp\beta)
=
\sup_{P^2=P=\transp{P},\ \tr(P)=k}\Wass_{2,P}(\alpha,\beta).
\]
\end{defn}
\(\gamma_k(M)=\sum_{\ell=1}^k\lambda_\ell(M),\)
\(\mathcal B_{\gamma_k}=\enscond{A}{0\preceq A\preceq\Id,\ \tr(A)\leq k}.\)
Thus $k=d$ gives $\gamma_d(M)=\tr(M)$ and recovers $\Wass_2$. The convex hull of rank-$k$ projectors is

\(\enscond{A}{0\preceq A\preceq\Id,\ \tr(A)=k},\)
and, since $M\succeq0$, the associated support function is the same Ky Fan gauge. Thus $\Wass_{\gamma_k}$ is the convexified spectral counterpart of $\SRW_{2,k}$, while $\SRW_{2,k}$ keeps the original non-convex rank constraint.

\[
\SRW_{2,k}(\alpha,\beta)
\leq
\Wass_{\gamma_k}(\alpha,\beta)
\leq
\Wass_2(\alpha,\beta),
\qquad
\sqrt{\frac{k}{d}}\,\Wass_2(\alpha,\beta)
\leq
\Wass_{\gamma_k}(\alpha,\beta),
\]
\[
\left\{\frac{k}{d}\Id\right\}
\cup
\enscond{P}{P^2=P=P^\top,\ \tr(P)=k}
\subset \mathcal B_{\gamma_k}
\subset
\enscond{A}{0\preceq A\preceq\Id}.
\]
For $k=1$, $\gamma_1(M)=\lambda_{\max}(M)$ and $\mathcal B_{\gamma_1}=\enscond{A\succeq0}{\tr(A)\leq1}$.

\subsection{Conditional Wasserstein Distances}
\label{sec-conditional-wasserstein-distances}
\begin{defn}[Conditional couplings and conditional OT]\label{def-conditional-ot}
Let $(S,\lambda)$ be a standard Borel probability space of conditions and let $(\Omega,\dist)$ be a Polish metric space. For $p\geq1$, denote by $\Pp_{p,\lambda}(S\times\Omega)$ the set of probability measures on $S\times\Omega$ whose first marginal is $\lambda$ and whose disintegration
\(\alpha(\d s,\d x)=\alpha_s(\d x)\lambda(\d s)\)
satisfies $\int_S\int_\Omega \dist(x,x_0)^p\,\d\alpha_s(x)\d\lambda(s)<+\infty$ for some, hence every, $x_0\in\Omega$.
For $\alpha,\beta\in\Pp_{p,\lambda}(S\times\Omega)$, a conditional coupling is a measure
\(\Pi(\d s,\d x,\d y) = \pi_s(\d x,\d y)\lambda(\d s)\)
such that $\pi_s\in\Couplings(\alpha_s,\beta_s)$ for $\lambda$-a.e. $s$. The set of such couplings is denoted $\Couplings_\lambda(\alpha,\beta)$.
Let $(s,x,y)\mapsto c_s(x,y)$ be jointly Borel, nonnegative, and lower semicontinuous in $(x,y)$ for every $s$. The conditional OT value is
\begin{equation}\label{eq-conditional-ot-general}
\begin{aligned}
\MK_c^\lambda(\alpha,\beta)
&\eqdef
\inf_{\Pi\in\Couplings_\lambda(\alpha,\beta)}
\int_S\int_{\Omega\times\Omega} c_s(x,y)\d\pi_s(x,y)\d\lambda(s) \\
&=
\int_S
\inf_{\pi_s\in\Couplings(\alpha_s,\beta_s)}
\int_{\Omega\times\Omega} c_s(x,y)\d\pi_s(x,y)
\d\lambda(s).
\end{aligned}
\end{equation}
\end{defn}
\begin{defn}[Conditional Wasserstein distance]\label{def-conditional-wasserstein-distance}
For $c_s(x,y)=\dist(x,y)^p$, define
\begin{equation}\label{eq-conditional-wasserstein-distance}
\Wass_{p,\lambda}(\alpha,\beta)
\eqdef
\left(
\int_S \Wass_p^p(\alpha_s,\beta_s)\d\lambda(s)
\right)^{1/p}.
\end{equation}
\end{defn}
\begin{prop}[Metric property]\label{prop-conditional-wasserstein-distance}
For $p\geq1$, $\Wass_{p,\lambda}$ is a distance on $\Pp_{p,\lambda}(S\times\Omega)$. Moreover, this metric space is Polish.
\end{prop}
\begin{proof}
Non-negativity and symmetry follow from the corresponding properties of $\Wass_p$ on each fiber. If $\Wass_{p,\lambda}(\alpha,\beta)=0$, then $\Wass_p(\alpha_s,\beta_s)=0$ for $\lambda$-a.e. $s$, hence $\alpha_s=\beta_s$ for $\lambda$-a.e. $s$ and the disintegrated measures $\alpha$ and $\beta$ coincide. For the triangle inequality, let $\gamma\in\Pp_{p,\lambda}(S\times\Omega)$. Since $\Wass_p(\alpha_s,\gamma_s)\leq\Wass_p(\alpha_s,\beta_s)+\Wass_p(\beta_s,\gamma_s)$ for a.e. $s$, Minkowski's inequality in $L^p(S,\lambda)$ gives
\(\Wass_{p,\lambda}(\alpha,\gamma) \leq \Wass_{p,\lambda}(\alpha,\beta)+\Wass_{p,\lambda}(\beta,\gamma).\)
To see completeness and separability, identify a measure with the $\lambda$-a.e. equivalence class of its measurable disintegration map $s\mapsto\alpha_s\in\Pp_p(\Omega)$. Formula~\eqref{eq-conditional-wasserstein-distance} is the corresponding $L^p(S,\lambda;\Pp_p(\Omega))$ metric. Since $(\Pp_p(\Omega),\Wass_p)$ is Polish, the standard $L^p$ space of measurable metric-valued maps is Polish as well.
\end{proof}
\begin{prop}[Conditional geodesics]\label{prop-conditional-wasserstein-geodesics}
Assume that $(\Omega,\dist)$ is a geodesic Polish space and let $\alpha_0,\alpha_1\in\Pp_{p,\lambda}(S\times\Omega)$. For $\lambda$-a.e. $s$, choose a constant-speed $\Wass_p$ geodesic $t\mapsto\alpha_{t,s}$ between $\alpha_{0,s}$ and $\alpha_{1,s}$, measurably in $s$, and set
\(\alpha_t(\d s,\d x)=\alpha_{t,s}(\d x)\lambda(\d s).\)
Then $t\mapsto\alpha_t$ is a constant-speed geodesic for $\Wass_{p,\lambda}$.
\end{prop}
\begin{proof}
For $0\leq r\leq t\leq1$, the fiberwise geodesic property gives \(\Wass_p(\alpha_{r,s},\alpha_{t,s}) = (t-r)\Wass_p(\alpha_{0,s},\alpha_{1,s}) \quad\text{for }\lambda\text{-a.e. }s.\)
Integrating the $p$th power over $S$ gives \(\Wass_{p,\lambda}(\alpha_r,\alpha_t) = (t-r)\Wass_{p,\lambda}(\alpha_0,\alpha_1).\)
Hence the conditional curve has constant speed and realizes the distance between its endpoints. In Euclidean space one may take, for each $s$, an optimal plan $\pi_s\in\Couplings(\alpha_{0,s},\alpha_{1,s})$ and set $\alpha_{t,s}=((1-t)\mathrm p_1+t\mathrm p_2)_\sharp\pi_s$, where $\mathrm p_1,\mathrm p_2$ are the two coordinate projections. On a general geodesic space, the same construction uses a measurable family of optimal dynamical plans. Such a family can be selected measurably in the Polish setting; when geodesics are non-unique, different selections can produce different conditional geodesics.
\end{proof}